========================================================
PROFIL J1 — STAGIAIRE, 5 MOIS D'ANCIENNETÉ
========================================================

Fonction : Stagiaire
Ancienneté : 5 mois
Missions : Propositions commerciales, chef de projet, support institutions européennes
Rapport à l'IA : Passionné, early adopter
Tic de langage : "Typiquement"
Contexte spécifique : A suivi l'évolution de Copilot PowerPoint de mars (inutilisable) à aujourd'hui (génération de slides complètes). Utilise Copilot officiellement et Claude officieusement (sans données sensibles).

--------------------------------------------------------
T0 — CADRAGE (parcours et missions)
--------------------------------------------------------

"Je suis arrivé en stage il y a cinq mois, un peu par hasard — je cherchais une expérience en conseil et Wavestone Luxembourg recrutait. Typiquement, je fais beaucoup de propositions commerciales et du support sur des missions pour des institutions européennes. La partie chef de projet aussi, donc je suis souvent en interface entre les seniors et les équipes techniques. La part d'écrit dans mon travail ? Énorme. Je dirais 70% de mon temps — Powerpoints, notes de cadrage, comptes-rendus de réunions. Le reste c'est de l'oral : des points avec le manager, des calls clients."

--------------------------------------------------------
T1 — USAGE RÉEL (Appropriation)
--------------------------------------------------------

Ouverture — dernière fois que vous avez utilisé un outil génératif :

"La dernière fois que j'ai utilisé un outil génératif au travail, c'était la semaine dernière sur une proposition commerciale. On avait un dossier de réponse à rendre et j'avais une dizaine de slides à structurer sur un sujet que je maîtrisais moyennement. Je suis allé sur Copilot directement dans PowerPoint — pas de copier-coller depuis ChatGPT — et je lui ai demandé de me générer une structure avec des titres et des sous-parties sur la base de trois documents qu'on avait déjà. Ça m'a pris cinq minutes alors que j'aurais passé une heure à tout relire pour faire le plan. Après j'ai repris chaque slide et j'ai réécrit les parties qui ne collaient pas au contexte [client]. Typiquement, le contenu généré est assez générique, mais ça te donne un squelette et tu n'as plus qu'à mettre la viande autour."

Relance — sur quelles tâches au quotidien ?

"Alors, tâches où je l'utilise tout le temps : la reformulation. Quand j'ai un texte un peu fouilli de notes de réunion, je le colle dans Copilot et je lui dis 'rends-moi ça structuré et fluide'. Là-dessus, c'est un tueur. Aussi pour les traductions — on bosse avec des institutions européennes donc j'ai souvent des documents en anglais ou en français à basculer de l'un à l'autre. Par contre, je ne l'utilise jamais pour les slides où il y a des données chiffrées confidentielles du client. Ça, c'est un dealbreaker."

Relance — comment as-tu appris à t'en servir ?

"Seul, clairement. Le cabinet n'a pas organisé de formation sur Copilot jusqu'à récemment — je crois qu'ils ont fait un webinaire le mois dernier, mais j'étais en mission. J'ai appris en trifouillant, en testant des prompts, en voyant ce qui marchait ou pas. Et surtout, j'ai suivi de très près les progrès : en mars, Copilot dans PowerPoint ne faisait quasi rien. Tu pouvais générer du texte brut sur un diapo mais sans mise en forme — c'était moche, inutilisable. En mai, on a eu des actions qui ne nécessitaient pas de styling, genre 'ajoute une transition' ou 'change la couleur du titre'. Et là, récemment, on peut générer des diapos entières avec une mise en forme correcte en un seul prompt. Typiquement, ça a changé la donne en trois mois. Mais j'ai dû apprendre par moi-même à chaque étape."

Relance — le temps gagné, tu l'as mis sur quoi ?

"La dernière fois, sur une proposition, j'ai gagné deux heures sur la structure et la rédaction initiale. Je les ai utilisées pour relire trois fois le document et vérifier que chaque recommandation était bien ancrée dans ce que le client avait dit en réunion. Donc je dirais que l'outil m'a permis de passer plus de temps sur la qualité de fond et moins sur la mise en forme et le squelette. C'est un peu contre-intuitif parce qu'on pense que l'IA va te faire gagner du temps pour en faire plus — en vrai, moi, elle m'a fait gagner du temps pour faire mieux sur l'essentiel."

Relance — une situation où ça t'a demandé plus d'effort que de le faire toi-même ?

"Ah oui, une fois. J'ai voulu lui faire générer un slide sur un sujet hyper technique lié à la réglementation bancaire. L'outil m'a sorti un texte vague, avec des formulations qui étaient presque justes mais pas tout à fait — genre des références à des directives qui n'existent pas. J'ai passé vingt minutes à lui demander de corriger, à reformuler mon prompt, à essayer de lui filer des sources. Au final, j'ai tout réécrit moi-même en quinze minutes. Typiquement, là, l'outil m'a coûté du temps au lieu de m'en faire gagner. Depuis, je sais : quand le sujet est trop pointu ou que je n'ai pas les sources exactes à lui donner, je ne perds pas mon temps. Je fais direct."

Relance — vous en parlez entre vous ?

"Oui, pas mal. Entre stagiaires et juniors, on se partage des astuces — des prompts qui marchent, des outils alternatifs. Mais avec les seniors, c'est plus discret. Typiquement, je ne vais pas dire à mon manager 'j'ai utilisé Copilot pour la structure de la proposition', même si elle le devine peut-être. C'est un peu un non-dit : on sait que tout le monde l'utilise, mais personne ne le crie sur les toits. Surtout les juniors — on a peur qu'on nous dise 'tu ne réfléchis plus par toi-même'."

--------------------------------------------------------
T2 — APPRENTISSAGE DU MÉTIER (Compétences)
--------------------------------------------------------

Ouverture — comment on apprend le métier aujourd'hui ?

"Alors, la façon dont j'apprends le métier… je pense que c'est assez différent de ce que vivait un stagiaire il y a cinq ans. Déjà, la partie rédactionnelle — les notes, les comptes rendus — je les fais beaucoup plus vite, mais je ne suis pas sûr d'apprendre la même chose. Je m'explique : quand tu rédiges un compte rendu à la main, tu es obligé de trier, de reformuler, de hiérarchiser les infos. Avec l'IA, je colle mes notes brutes et je lui dis 'structure-moi ça'. Je gagne du temps, mais je n'exerce pas le muscle de la synthèse. Typiquement, je sens que je suis plus un 'relecteur' qu'un 'rédacteur' maintenant. Je vérifie ce que l'IA a fait, mais je n'ai pas passé les vingt minutes à le construire moi-même."

Relance — par quelles tâches on apprend réellement le métier ?

"Pour moi, ce qui m'a le plus appris, c'est la relecture par le manager. La première fois que j'ai fait une proposition commerciale, j'ai passé deux jours à la construire, puis mon manager m'a renvoyé un document rempli de commentaires — des titres pas clairs, une structure bancale, des arguments mal placés. J'ai tout repris en comprenant *pourquoi*. C'était dur mais formateur. Maintenant, avec Copilot, j'arrive avec une première version assez propre très vite. Mon manager passe moins de temps sur la forme et plus sur le fond. Du coup, ses commentaires sont plus intéressants — il me dit 'ton raisonnement est léger sur tel point' plutôt que 'tes titres ne sont pas alignés'. Donc paradoxalement, l'IA a amélioré la qualité de la relecture senior. Mais je me demande si, à force, je vais perdre la capacité à faire une première version toute seule. Typiquement, c'est ma crainte."

Relance — les tâches ingrates, les CV par exemple, sont-elles encore formatrices ?

"Ah, les CV. C'est un bon exemple. Je dois souvent adapter des CV existants à des postes spécifiques — choisir les bonnes expériences, développer certaines missions plutôt que d'autres, reformuler pour coller à une offre. C'est un exercice qui demande de la nuance, de comprendre ce qui est pertinent. Au début, je le faisais à la main. Puis j'ai essayé avec Copilot : je lui filais le CV original, la description du poste, et je lui demandais de me proposer une version adaptée. Ça marchait plutôt bien. Mais une fois, et c'est là que ça a dérapé, l'IA est partie en cacahuète. Elle m'a sorti un CV avec des projets qui n'existaient pas — ou qui appartenaient à d'autres personnes, on n'a jamais su. Des trucs du genre 'chef de projet sur la refonte de X' alors que la personne n'avait jamais fait ça. Si je n'avais pas vérifié ligne par ligne avec le profil original, on aurait envoyé un CV totalement mensonger. Typiquement, là, je me suis dit : 'OK, l'outil est puissant mais il peut inventer n'importe quoi.' Depuis, je l'utilise pour gagner du temps sur la reformulation, mais je ne lui fais plus jamais confiance sur les faits."

--------------------------------------------------------
T3 — CLIENT ET VALEUR (Légitimité de l'expert)
--------------------------------------------------------

Ouverture — la confiance du client repose sur quoi ?

"La confiance du client, je pense que ça repose sur deux choses. La première, c'est notre capacité à parler leur langage — à comprendre leur métier, leurs contraintes. La deuxième, c'est notre responsabilité : si on leur donne un livrable, ils doivent pouvoir se dire 'ces gens-là ont vérifié, ils sont engagés'. Typiquement, un client ne va pas te demander si tu as utilisé l'IA ou pas, il va te demander si ce que tu lui présentes est vrai. Donc la confiance, elle est dans le fait qu'on a mis notre nom sur le document."

Relance — diriez-vous à un client qu'une partie a été produite par l'IA ?

"Honnêtement ? Non. Pas en l'état. Je ne mentirais pas si on me demande directement, mais je ne vais pas le mettre en avant. Il y a un vrai risque que le client se dise 'pourquoi je paie un cabinet si c'est un algorithme qui fait le boulot ?'. Donc je préfère ne pas en parler. Par contre, je pense que les clients le savent ou le soupçonnent. Certains nous ont déjà demandé, en réunion, si on utilisait des outils d'IA. On a éludé : 'on utilise toutes les technologies qui nous aident à être plus pertinents'. Mais on n'a jamais dit 'oui, on génère des slides avec Copilot'. Typiquement, c'est un sujet tabou."

Relance — si un livrable était perçu comme moins crédible, pourquoi ?

"Le problème, c'est quand l'IA fait des erreurs factuelles. Si un client repère un chiffre faux, il va remettre en cause tout le reste. Là, l'outil n'a pas de 'statut' qui le protège — ce n'est pas une personne, donc on ne peut pas dire 'c'est une erreur humaine'. L'erreur humaine est excusable, l'erreur machine est perçue comme une faute professionnelle. Donc je suis très prudent sur tout ce qui est chiffré ou factuel. Je vérifie tout. Typiquement, si je génère un slide avec des données, je repasse derrière avec les sources originales."

Relance — dans dix ans, sur quoi un client acceptera-t-il de payer un cabinet ?

"Je pense que le client paiera pour notre capacité à prendre des risques et à dire des choses qu'il n'a pas envie d'entendre. L'IA peut faire de l'analyse, de la synthèse, de la mise en forme. Elle ne peut pas prendre position. Elle ne peut pas dire au client 'vous vous trompez sur ce point' ou 'il faut que vous arrêtiez tel projet'. Le cabinet, il est payé pour le jugement. Typiquement, le jour où l'IA saura faire ça, le métier sera vraiment différent. Mais pour l'instant, la valeur ajoutée du consultant, c'est de dire des vérités inconfortables."

--------------------------------------------------------
T4 — GOUVERNANCE ET RISQUES
--------------------------------------------------------

Ouverture — il existe des règles sur l'usage des outils ici ?

"Alors, des règles... officiellement, oui. On a eu une session de sensibilisation sur l'IA générative il y a environ deux mois. On nous a dit : Copilot est l'outil approuvé par le cabinet, il est intégré à la suite Microsoft, donc vous pouvez l'utiliser dans le cadre de vos missions, mais attention aux données confidentielles — il faut anonymiser, ne pas coller des documents sensibles en entier. Concernant Claude, ChatGPT ou autre... c'est le flou total. On nous a dit 'faites attention' mais il n'y a pas de règle écrite qui interdit ou autorise clairement. Donc typiquement, je l'utilise officieusement, en faisant très attention à ce que je mets dedans. Mais je ne suis pas sûr que tout le monde soit aussi prudent."

Relance — avant qu'un livrable ne parte chez le client, il se passe quoi ?

"Il y a un processus de revue. En général, je produis une version préliminaire, je l'envoie à mon manager, il la réécrit ou met des commentaires, je corrige, puis on la renvoie pour validation finale. C'est cette relecture qui fait office de contrôle. Mais il n'y a pas de 'check IA' spécifique. Personne ne me demande 'est-ce que cette partie a été générée par l'outil ?' On me demande juste de justifier mes choix, mes chiffres, mes sources. Donc le contrôle est sur le fond, pas sur le processus. Je pense que certains managers soupçonnent qu'on utilise l'IA, mais ils ne posent pas la question directement."

Relance — l'histoire du CV où l'IA a inventé des projets, qu'est-ce qui s'est passé après ?

"Je l'ai dit à un collègue stagiaire, on en a rigolé, mais je ne l'ai pas dit à mon manager. J'avais un peu honte — je me suis dit qu'elle allait penser que je ne faisais pas attention ou que j'étais trop naïf de faire confiance à l'outil. Du coup, j'ai repris le CV, j'ai tout réécrit à la main, et je l'ai envoyé comme si de rien n'était. Mais typiquement, j'ai eu une vraie frayeur. Si j'avais envoyé ce CV sans vérifier, le client aurait reçu un profil avec des mensonges. Et là, c'est moi qui aurais été responsable, pas l'IA. Personne ne vient vérifier ce qui sort de l'outil à ma place, donc j'ai compris que la responsabilité, elle est entièrement sur moi."

Relance — si vous écriviez la règle vous-même, vous diriez quoi ?

"Je dirais : 'Vous pouvez utiliser l'IA pour tout ce qui est rédactionnel, structure, synthèse. Mais vous devez être capable de justifier chaque affirmation factuelle. Si vous ne pouvez pas la justifier, vous ne l'utilisez pas.' Je ne mettrais pas d'interdiction stricte parce que l'outil est trop utile. Mais je mettrais une obligation de vérification systématique, surtout sur les données, les références, les noms de projets. Typiquement, ce qui m'a sauvé sur le CV, c'est que j'avais gardé le profil original à côté pour comparer. Si j'avais juste envoyé le résultat brut, on aurait eu un vrai problème. Donc la règle, ce serait : 'toujours garder une source de vérité pour vérifier'."

--------------------------------------------------------
T5 — CLÔTURE
--------------------------------------------------------

Relance — à la place de la direction, quelle serait votre première décision ?

"Ma première décision, ce serait de clarifier le statut des outils non-officiels. Soit on interdit Claude et ChatGPT officiellement, soit on les approuve avec des garde-fous. Mais le flou actuel est dangereux — certains juniors s'en servent sans se poser de questions, d'autres n'osent pas du tout, et tout le monde fait un peu ce qu'il veut. Je formaliserais un 'usage autorisé' et un 'usage interdit' clairs. Et j'ajouterais une formation obligatoire — pas sur le prompt engineering, mais sur la détection des hallucinations. Parce que le vrai risque, ce n'est pas que l'outil fasse des erreurs, c'est que les consultants ne les voient pas."

Relance — un aspect important qu'on n'a pas abordé ?

"Le temps psychologique. Parce que l'outil te fait gagner du temps, mais il crée aussi une pression. Typiquement, si je sais que je peux générer une slide en 30 secondes, je suis moins enclin à passer du temps à réfléchir sur le fond avant de commencer. Je lance l'outil, il me donne un truc, je le modifie. Ça va plus vite, mais j'ai parfois l'impression que je ne maîtrise plus le sujet comme je le maîtrisais quand je faisais tout à la main. C'est un peu... une illusion de compétence. Je me demande si, dans six mois ou un an, je saurai encore faire une bonne structure tout seul sans l'aide de l'IA."

Relance — quelqu'un dont le point de vue serait particulièrement utile ?

"Un manager qui est là depuis dix ans. Pour voir comment il vit le fait que les juniors arrivent avec des livrables de plus en plus 'propres' mais peut-être moins réfléchis. Parce que c'est lui qui valide, lui qui signe. Je pense que son regard est différent du mien — moi je vois l'outil comme un accélérateur, lui pourrait le voir comme un risque. Typiquement, j'aimerais savoir s'il fait la différence entre un livrable fait par un junior compétent et un livrable fait par un junior qui a utilisé l'IA."

========================================================
FIN DU VERBATIM J1
========================================================